\section{A nonnegative two-sided bar still introduces a vertical tail}

\paragraph{ARCH-136-RELATIVE-BAR-OBSTRUCTION.}
This tests the common-target route in cote 136, physical pp.~117,119,121,
archive footers 116,118,120, supplied scans 117,119,121.
It does not refute the existence of another common target.

There exists a simplicial strict PD differential \(S=\Gamma_{\mathbb Z}(T)\)-algebra
\(H\), with two strict maps \(H\to B\), such that:
\begin{enumerate}
\item \(H_0=S\), \(H_1=B_1\), and every horizontal row of \(H\) is
simplicially homotopy equivalent to the corresponding row of \(B\);
\item every simplex has coefficient \(H^0=S\) and vanishing vertical
\(H^1,H^2,H^3\), in every weight;
\item the full differential subalgebra generated by the zero forms of
\(H\) is \(B\), and \(H\) is free as a graded module over \(B\);
\item the normalized relative two-sided bar of the two augmentations,
formed over this zero-form subalgebra, is nonnegative in total cochain
degree and has finite bar length in every weight;
\item nevertheless this bar has nonzero vertical \(H^3\) at simplex two
in every weight \(w\geq2m\), for one fixed positive integer \(m\).
\end{enumerate}
Thus removing zero-form bar letters and the indicated freeness problems
does not make this particular bar construction preserve vertical admissibility.

\begin{proof}
Let
\[
 U=\operatorname{DK}\bigl(\mathbb Z[3]\xrightarrow{1}\mathbb Z[2]\bigr).
\]
This is simplicially contractible, is free abelian in every degree, and
has \(U_0=U_1=0\), \(U_2=\mathbb Z\).
Let \(F(U)\) be the free strict graded PD algebra over \(S\) on copies
\(a(U)\) and \(b(U)\) of \(U\), in cochain degrees one and two respectively,
with \(da(u)=b(u)\), \(db(u)=0\). Give all these generators weight \(m>0\).
Define
\[
                       H=B\otimes_S F(U).
\]
Both maps to \(B\) are the augmentation of \(F(U)\).

A simplicial linear contraction of \(U\) consists of a family of linear
maps satisfying the simplicial homotopy identities. Apply the free PD
functor to those maps. Since a functor preserves each identity, this is
an algebra homotopy from \(F(U)\) to its coefficient augmentation.
Tensor the homotopy with the degeneracies of \(B\). This gives the
asserted horizontal homotopy equivalence, preserving form and weight.
It also proves the first-skeleton assertions.

At any simplex, choosing a basis of \(U_n\) writes \(F(U)_n\) as a finite
tensor of the odd disks
\[
 D=S\otimes\Lambda(a)\otimes\Gamma(b),\qquad da=b.
\]
A single disk splits, as a cochain complex, into its coefficient copy,
the contractible pair \(Sa\to Sb\), and the pairs
\[
 S\,a\gamma_{r-1}(b)\xrightarrow{r}S\,\gamma_r(b),
                       \qquad r\geq2.
\]
The latter start in form degree \(2r-1\geq3\), and the first such
differential is multiplication by two, which is injective in \(S\).
A tensor involving at least two of these noncontractible summands starts
in degree at least six; all summands involving a contractible factor
contract. Hence a finite tensor has \(H^1=H^2=H^3=0\), and coefficient
\(H^0=S\). The \(S\)-linear contraction of \(B_n\) preserves this conclusion.
All assertions are weightwise and valid for every simplex.

The positive generators do not change zero forms, so the zero-form
differential subalgebra is \(B\). The relative normalized two-sided bar
therefore is
\[
 E_n=B_n\otimes_S\overline{\operatorname{Bar}}_S(F(U)_n),
\]
where the reduced bar includes its length-zero coefficient term, each
letter is in the augmentation ideal, and suspension lowers its cochain
degree by one. Every letter has cochain degree at least one and weight
at least \(m\). The total is consequently nonnegative, and in weight
\(w\) its bar length is at most \(\lfloor w/m\rfloor\).
The modules involved are free over the stated base.

At simplex two there is exactly one odd disk. Use the following
cohomological bar convention. For homogeneous letters \(x,y\), put
\[
 d_1[x]=-[dx],\qquad d_2[x|y]=(-1)^{|x|}[xy],
\]
and extend the two operations as coderivations with signs determined
by the suspended degree \(|x|-1\). The internal differential and
multiplication satisfy \(d^2=0\).

Give \(a,b\) auxiliary degree one and all coefficients auxiliary degree
zero. The bar differential preserves this auxiliary degree. In auxiliary
degree two its basis and total cochain degrees are
\[
\begin{array}{c|c}
0&X=[a|a]\\
1&Y=[a|b],\quad Z=[b|a]\\
2&W=[b|b],\quad V=[ab]\\
3&L=[\gamma_2(b)] .
\end{array}
\]
Its differential is exactly
\[
\begin{gathered}
 dX=-Y-Z,\qquad dY=-W-V,\qquad dZ=W+V,\\
 dW=2L,\qquad dV=-2L,\qquad dL=0 .
\end{gathered}
\]
The first two successive images are primitive kernels. The final image
is \(2S L\). Thus this direct summand has cohomology \(S/2S\) in degree
three, shifted by weight \(2m\), and zero in its other degrees.
Tensoring with the explicit weight-preserving \(S\)-linear contraction
of \(B_2\) to \(S\) leaves this cohomology unchanged. It follows that
\[
              H^3(E_2^{w,*})\supseteq S^{w-2m}/2S^{w-2m}
                        \cong\mathbb Z/2\qquad(w\geq2m).
\]
This is the required fixed-form-degree vertical tail.
\end{proof}

\paragraph{Scope of the bar claim.}
The obstruction is in the explicitly displayed underlying normalized
bar complex. It is therefore present even if all the desired shuffle
products and strict PD operations have been supplied. The proof does
not need to infer a positive cochain model structure from a homological
bar theorem. It does not say that the ordinary underived amalgamated
tensor, a quotient of the bar, or a different relative base must fail.
In this example the two original target theories are both \(B\), so
they already have the identity common target.

